% =============================================================================
%  Datei-Hinweise & Konfiguration: siehe config.md
%  -> Abschnitt "verzeichnisse/formelverzeichnis-defs.tex"
%
%  Formelzeichen des Formelverzeichnisses (Leitfaden Kap. 5.1 / 5.8). Eingebunden via
%  \loadglsentries in main.tex; gedruckt in verzeichnisse/formelzeichen.tex.
%  Der sort-Schluessel ist bewusst als Klartext gesetzt, da die Namen
%  Mathematik enthalten (noidx-Sortierung).
% =============================================================================

\glsxtrnewsymbol[sort={B}, description={Brier-Score: mittlere quadratische
    Abweichung prognostizierter Wahrscheinlichkeiten vom eingetretenen
    Ergebnis}]{symBrier}{\ensuremath{B}}
\glsxtrnewsymbol[sort={C(S)}, description={Round-Trip-Transaktionskosten bei
    Kapitaleinsatz $S$ (Fixgebühren zuzüglich Spread-Anteil)}]{symCS}{\ensuremath{C(S)}}
\glsxtrnewsymbol[sort={K}, description={Basiskapital der Strategie
    (2.000\,\euro{}); zugleich harte Obergrenze des Kapitaleinsatzes}]{symK}{\ensuremath{K}}
\glsxtrnewsymbol[sort={n}, description={Fensterlänge des gleitenden
    Durchschnitts in Handelstagen}]{symN}{\ensuremath{n}}
\glsxtrnewsymbol[sort={N}, description={Anzahl der bewerteten Prognosetage der
    Walk-Forward-Prüfung}]{symNBig}{\ensuremath{N}}
\glsxtrnewsymbol[sort={ot}, description={eingetretenes Ergebnis am Handelstag
    $t$ (1, falls die Folgetagsrendite in Signalrichtung positiv ausfiel,
    sonst 0)}]{symOt}{\ensuremath{o_t}}
\glsxtrnewsymbol[sort={pt}, description={prognostizierte Wahrscheinlichkeit
    einer positiven Folgetagsrendite in Signalrichtung}]{symPt2}{\ensuremath{p_t}}
\glsxtrnewsymbol[sort={p*(S)}, description={Break-even-Trefferquote: für einen
    positiven Erwartungswert erforderliche Trefferquote bei Kapitaleinsatz
    $S$}]{symPstar}{\ensuremath{p^{*}(S)}}
\glsxtrnewsymbol[sort={P(t)}, description={Tagesschlusskurs des S\&P~500 am
    Handelstag $t$}]{symPt}{\ensuremath{P(t)}}
\glsxtrnewsymbol[sort={R(t)}, description={Handelsregime der
    Kreuzungsstrategie am Handelstag $t$ (LONG oder SHORT)}]{symRt}{\ensuremath{R(t)}}
\glsxtrnewsymbol[sort={R'(t)}, description={Handelsregime der
    Trendfolge-Alternative \emph{smatrend} (LONG oder CASH)}]{symRpt}{\ensuremath{R'(t)}}
\glsxtrnewsymbol[sort={S(t)}, description={volatilitätsabhängig skalierter
    Kapitaleinsatz am Handelstag $t$}]{symSt}{\ensuremath{S(t)}}
\glsxtrnewsymbol[sort={Smin}, description={Mindestpositionsgröße
    (1.000\,\euro{})}]{symSmin}{\ensuremath{S_{\min}}}
\glsxtrnewsymbol[sort={SMAn(t)}, description={einfacher gleitender Durchschnitt
    der Tagesschlusskurse über $n$ Handelstage}]{symSMA}{\ensuremath{\mathrm{SMA}_n(t)}}
\glsxtrnewsymbol[sort={sigma*}, description={Zielvolatilität der
    Positionsgrößensteuerung (1{,}0\,\% pro Tag)}]{symSigStar}{\ensuremath{\sigma^{*}}}
\glsxtrnewsymbol[sort={sigma20}, description={realisierte 20-Tage-Volatilität
    der Index-Tagesrenditen}]{symSig20}{\ensuremath{\sigma_{20}(t)}}
\glsxtrnewsymbol[sort={t}, description={Index des Handelstags}]{symT}{\ensuremath{t}}

%  =============================================================================
%  Ende
%  =============================================================================
